\scanpage{19}
% L1–3
On vérifie que cette application est fonctorielle en $Z$, donc définit $j_{u_0}\colon u_0\longrightarrow\alpha\beta(u_0)$, et que cette dernière est fonctorielle en $u_0$, d'où un hom de foncteurs

% L4
\[
\boxed{j\colon\mathrm{id}_{\underline{\mathrm{Hom}}(C,E)}\longrightarrow\alpha\beta}
\]

% L5–6
Signalons tout de suite que pour \deleted{\unsure{$Z$|que}} $u_0$ donné, $j_{u_0}\colon u_0\longrightarrow\alpha\beta(u_0)$ est un isom. sss \deleted{\unsure{on}} $\forall Z\in\mathrm{Ob}\,C$,

% L7
\[
u_0(z)\to\varinjlim_{C/\varepsilon_C(z)}u_0(X)\qquad\text{est un isom,}
\]

% L8–9
ou encore sss pour tout tel $Z$ et $\xi\in\mathrm{Ob}\,E$, la flèche déduite de la précédente en appliquant $\mathrm{Hom}(-,\xi)$ est bijective,

% L10
\[
\text{i.e.}\qquad\varprojlim_{C/\varepsilon_C(z)}\mathrm{Hom}_{\unsure{\iota|1|{}}}(u_0(X),\xi)\longrightarrow\mathrm{Hom}(u_0(z),\xi)\qquad\text{bijective.}
\]

% L11
Or soit \deleted{\unsure{pour $\xi\in\mathrm{Ob}\,E$}},

% L12
\[
u_0^{\unsure{\bullet|\cdot|\prime}}\colon E\longrightarrow\widehat C
\]

% L13
le foncteur défini par

% L14
\[
u_0^\bullet(\xi)(Y)=\mathrm{Hom}(u_0(Y),\xi)
\]

% L15
de sorte que l'application précédente devient, posant $P=u_0^\bullet(\xi)$ :

% L16
\[
\varprojlim_{C/\varepsilon_C(z)}P(X)\longrightarrow P(Z)\qquad\text{un isom.}
\]

% L17–19
\deleted{\unsure{Nous proposons}} Cette condition est vide si $\varepsilon_C$ est pl. fidèle i.e. la top. de $C$ est moins fine que la top. canonique (On sait que $C/\varepsilon_C(z)\simeq C/z$ admet $Z$ comme objet final).

% L20–24
On va prouver que $i,j$ sont des données d'adjonction associées pour les foncteurs $\alpha$ et $\beta$, [\unsure{Celà|Cela|Ce} qui prouve que $\alpha$ et $\beta$ \deleted{établissent} \inserted{induisent} des équivalences entre les sous-catégories \inserted{str. pleines} de $\underline{\mathrm{Hom}}\,(\widetilde C,E)$ et $\underline{\mathrm{Hom}}\,(C,E)$ qu'on a décrit \unsure{auparavant|au par. n$^\circ$ 7|au préal}. On \deleted{\unsure{en conclut}} \inserted{en conclut} \deleted{\unsure{a établi}} que $\alpha$ induit un foncteur \uline{pl. fidèle}

% L25
\[
\underline{\mathrm{Hom}}'(\widetilde C,E)\longrightarrow\underline{\mathrm{Hom}}\,(C,E)
\]

% L26–28
dont l'image essentielle est formée des \deleted{\unsure{$Z$}} $u_0\colon C\longrightarrow E$ satisfaisant la condition ci-dessus, et tels que $\beta(u_0)\in\underline{\mathrm{Hom}}'(\widetilde C,E)$ -- condition qu'on remplacera par une autre plus simple]

% foot
% archive sheet number: 18
